\documentclass[11pt]{scrartcl}

\usepackage[stixtwo]{fontsetup}
\usepackage{amsmath}

\begin{document}
A függvény folytonos függvények kompozíciója, így a nevező zérushelyeiben,
tehát az $x = -2$ és az $x = 5$ pontban van, valamint az $x = 0$ pontban lehet
szakadása, máshol mindenütt folytonos. ($x = 4$ esetén nincs
baj, hiszen ott máshogy van értelmezve a függvény.)
\[
	\lim_{x \to -2 \pm 0} f(x)
	= \lim_{x \to -2 \pm 0} \frac{\sin(x + 2)}{x + 2} \cdot \frac{1}{x - 4}
	= 1 \cdot \left( -\frac{1}{6} \right),
\]
vagyis ebben a pontban a függvénynek megszüntethető szakadása van.

\[
	\lim_{x \to 5 \pm 0} f(x)
	= \lim_{x \to 5 \pm 0} \frac{3x^2}{x - 5}
	= \pm \infty,
\]
vagyis ebben a pontban a függvénynek másodfajú szakadása van.
\[
	\lim_{x \to 0 - 0} f(x)
	= f(0)
	= \frac{-\sin 2}{8}
	\ne 0
	= \frac{3 \cdot 0^2}{0 - 5}
	= \lim_{x \to 0 + 0} \frac{3x^2}{x - 5},
\]
így az $x = 0$ pontban véges ugrása van a függvénynek.
\end{document}
